\cotefile{8}
\scanpage{17}
% sheet: cancelled_draft
% archive sheet number: 16
% Three long oblique pencil strokes cross the two sections.
% @ text / top, item 3
\deleted{-- \textcircled{3} Indétermination dans ces données, pour $(\deleted{\unsure{G}}M,T)$ (i.e. $G$) fixée.}

% @ text / upper third, body of item 3
\deleted{La catégorie des données \deleted{$\unsure{T}$\illegible{1}} en question est \deleted{\illegible{1}} ``un torseur'' sous la catégorie des torseurs sous $\mu_2$ sur $A'$, donc l'ensemble des classes d'isomorphie de ces structures est \inserted{un torseur sous} $H^1(A^{\deleted{\unsure{*}}},\mu_2)\simeq\prod_p Q_p^*/Q_p^{*2}$, \inserted{\unsure{pr. direct|pr. dir.} local rel. à \unsure{l'un des ss-gpes|l'un des sous-gpes|l'un des deux-gpes} $\mathbb Z_p^*/\mathbb Z_p^{*2}$} l'ensemble des autom. d'un objet de la catégorie est $H^0(A',\mu_2)=\prod_p(\mathbb Z/2\mathbb Z)$}
% slanted
\margin{modulo \unsure{mes|une} \unsure{légers renforcements de la|ligne renforcée de la} donnée à l'\unsure{$\infty$}}

% @ text / middle, item 4
\deleted{\textcircled{4} Rigidité de $M$ pour les structures envisagées.}

% @ text / middle, before equivalence
\deleted{Soit $(M',T')$ muni des \deleted{\unsure{mêmes}} données \inserted{suppl.} que $(M,T)$. On se demande de déterminer la catégorie des équivalences}

% @ display / middle
\deleted{\[
(M,T,\overbrace{T_{A'},u_{A'}}^{t_{A'}},t_\infty,t_p)\longrightarrow (M',T'\ \text{etc}\qquad)
\]}

% @ text / middle to lower third
\deleted{\deleted{\unsure{Elles|Ils|Elle}} \inserted{Cette catégorie est vide, \uuline{ou}} forme un torseur sous la catégorie des torseurs sous $\mu_2$ sur $Q$, \uline{munis} d'une trivialisation \deleted{\unsure{de}\illegible{1}} au dessus \deleted{\unsure{de}} $A$. Cette dernière catégorie a comme \deleted{\unsure{ensemble}} \inserted{groupe} de \deleted{\unsure{d'id}} classes d'objet à isom près le groupe}

% @ display / lower middle
\deleted{\[
H^{1}(Q\ \mathrm{mod}\ A,\mu_2)\simeq {}_2(\unsure{J|\mathcal I}_Q)\simeq\bigl(\prod_p\mathbb Z/2\mathbb Z\bigr)/(\pm1)
\]}
\margin{qui a la puiss. du continu $\bigr]$}

% @ text / lower third, rigidity
\deleted{d'autre part cette catégorie est rigide $(H^0(Q\ \mathrm{mod}\ A,\mu_2)=0)$.}

% @ text / foot
\deleted{Comme l'image de $H^1(Q\ \mathrm{mod}\ A,\mu_2)$ dans $H^{1}(Q,\mu_2)$ est nulle, on voit que \uline{l'équivalence induite entre} $(M,T)$ \uline{et} $(M',T')$ \uline{est déterminée à isom. unique près}. Donc on peut identifier sans dommages $(M,T)$ et $(M',T')$}
